\chapter{Bernalar dalam Ketidakpastian}
\label{ch:prob}

Kuliah model probabilistik saya ikuti pada musim dingin 2025/26. Slidenya banyak bersandar pada
buku \citet{koller2009} dan \citet{russell2021}. Bab ini mengikuti urutan kuliahnya: representasi,
inferensi, pembelajaran, lalu model yang berjalan dalam waktu.

Bab-bab sebelumnya kebanyakan membahas fungsi yang memetakan input ke output. Bab ini membahas cara
lain memandang dunia: sebagai kumpulan variabel acak yang saling bergantung. Kita tidak lagi
bertanya ``apa jawabannya?'', tetapi ``seberapa yakin kita pada setiap kemungkinan jawaban, dan
bagaimana keyakinan itu berubah ketika ada bukti baru?''.

\section{Mengapa butuh struktur}

Distribusi gabungan atas $n$ variabel biner butuh $2^n-1$ angka. Untuk dua puluh variabel, itu
sekitar sejuta angka, dan untuk tiga puluh variabel, satu miliar. Tidak ada yang bisa menyimpan,
apalagi mempelajari, tabel sebesar itu. Jalan keluarnya adalah \istilah{conditional independence}.
Variabel $X$ dan $Y$ independen bersyarat terhadap $Z$ kalau, setelah $Z$ diketahui, mengetahui $Y$
tidak lagi menambah informasi tentang $X$: $p(X\mid Y,Z)=p(X\mid Z)$. Petir dan suara guntur sangat
berkorelasi, tetapi keduanya disebabkan oleh peristiwa yang sama. Kalau kita sudah tahu ada sambaran
listrik di awan, mendengar guntur tidak menambah informasi tentang kilatnya.

Aturan rantai probabilitas selalu berlaku: $p(x_1,\dots,x_n)=\prod_ip(x_i\mid x_1,\dots,x_{i-1})$.
Independensi bersyarat memangkas setiap faktor sehingga hanya bergantung pada beberapa variabel.

\section{Bayesian network}

\istilah{Bayesian network} adalah graf berarah tanpa siklus yang setiap simpulnya adalah variabel acak.
Setiap variabel diberi distribusi bersyarat terhadap induk-induknya, dan distribusi gabungannya
adalah hasil kali semuanya:
\begin{equation}
p(x_1,\dots,x_n)=\prod_{i=1}^n p\big(x_i\mid\mathrm{pa}(x_i)\big).
\end{equation}
Kalau setiap variabel biner punya paling banyak $k$ induk, jaringan butuh paling banyak $n\,2^k$
angka. Dengan $n=20$ dan $k=3$, itu 160 angka, bukan sejuta.

Contoh yang dipakai banyak buku teks, termasuk \citet{koller2009}, adalah seorang mahasiswa. Nilai
ujiannya, $G$, bergantung pada tingkat kesulitan mata kuliah, $D$, dan kecerdasan mahasiswa, $I$.
Kecerdasan juga memengaruhi skor tes masuk, $S$, dan nilai memengaruhi kualitas surat rekomendasi,
$L$. Graf ini memungkinkan beberapa pola penalaran. Penalaran kausal berjalan searah panah: mahasiswa
cerdas cenderung mendapat nilai baik. Penalaran evidensial berjalan melawan panah: nilai buruk
menurunkan keyakinan bahwa mahasiswanya cerdas. Pola yang paling menarik adalah penalaran
antarsebab. Nilai buruk menaikkan keyakinan bahwa mata kuliahnya sulit dan menurunkan keyakinan
bahwa mahasiswanya cerdas. Kalau kemudian kita tahu mata kuliahnya memang sulit, nilai buruk
itu sudah ``terjelaskan'', dan keyakinan tentang kecerdasan mahasiswa naik kembali. Pola ini
disebut \istilah{explaining away}. Bayesian network menangani semua pola ini dengan satu algoritme
inferensi yang sama, tanpa aturan khusus untuk masing-masing (\cref{fig:mahasiswa}).

\begin{figure}[ht]
\centering
\begin{tikzpicture}[x=1.8cm,y=1.2cm]
  \node[bulat,minimum size=9mm] (D) at (0,1) {$D$};
  \node[bulat,minimum size=9mm] (I) at (2,1) {$I$};
  \node[bulat,minimum size=9mm] (G) at (1,0) {$G$};
  \node[bulat,minimum size=9mm] (S) at (3,0) {$S$};
  \node[bulat,minimum size=9mm] (L) at (1,-1) {$L$};
  \draw[panah] (D) -- (G);
  \draw[panah] (I) -- (G);
  \draw[panah] (I) -- (S);
  \draw[panah] (G) -- (L);
  \node[label,left] at (D.west) {kesulitan};
  \node[label,right] at (I.east) {kecerdasan};
  \node[label,left] at (G.west) {nilai};
  \node[label,right] at (S.east) {skor tes};
  \node[label,left] at (L.west) {surat rekomendasi};
\end{tikzpicture}
\caption{Jaringan mahasiswa. Distribusi gabungannya
$p(D,I,G,S,L)=p(D)\,p(I)\,p(G\mid D,I)\,p(S\mid I)\,p(L\mid G)$.}
\label{fig:mahasiswa}
\end{figure}

Explaining away paling mudah dilihat dengan angka pada contoh yang lebih kecil. Rumput basah ($W$) bisa disebabkan
hujan ($R$) atau penyiram ($S$), yang saling independen dengan $p(R)=0{,}2$ dan $p(S)=0{,}1$, dan rumput basah tepat
ketika salah satunya terjadi. Peluang rumput basah adalah $1-0{,}8\cdot0{,}9=0{,}28$. Begitu kita melihat rumput basah,
peluang hujan naik dari $0{,}2$ menjadi $p(R\mid W)=0{,}2/0{,}28\approx0{,}71$. Lalu kita melihat penyiram sedang menyala.
Penyiram sudah cukup menjelaskan rumput basah, sehingga rumput basah tidak lagi memberi informasi tentang hujan, dan
$p(R\mid W,S)$ kembali ke $0{,}2$. Dua sebab yang semula independen menjadi saling bergantung begitu akibat bersamanya
diamati.

\subsection{Membaca independensi dari graf}

Kapan dua variabel independen bersyarat bisa dibaca langsung dari bentuk graf. Ada tiga pola dasar
untuk jalur $X-Z-Y$. Pada rantai $X\to Z\to Y$ dan pada garpu $X\leftarrow Z\to Y$, jalur informasi
terbuka selama $Z$ tidak diamati dan tertutup ketika $Z$ diamati. Pola ketiga, \istilah{v-structure}
$X\to Z\leftarrow Y$, bekerja terbalik. Kalau $Z$ tidak diamati, jalurnya tertutup: kesulitan mata
kuliah dan kecerdasan mahasiswa tidak saling memengaruhi. Kalau $Z$, atau salah satu turunannya,
diamati, jalur terbuka, dan itulah sebabnya explaining away terjadi. Dua variabel disebut
\istilah{d-separated} oleh sekumpulan pengamatan kalau semua jalur di antara keduanya tertutup, dan
dalam hal itu mereka pasti independen bersyarat.

Satu himpunan yang istimewa adalah \istilah{Markov blanket} sebuah variabel: induknya, anaknya, dan
induk lain dari anak-anaknya. Kalau Markov blanket diketahui, variabel itu independen dari seluruh
sisa jaringan. Sifat ini nanti menjadi kunci Gibbs sampling.

\section{Inferensi eksak}

Inferensi berarti menghitung distribusi sebuah variabel kueri setelah sebagian variabel lain
diamati, misalnya $p(I\mid G=\text{buruk})$. Cara naif adalah menjumlahkan distribusi gabungan
atas semua variabel lain, dengan biaya eksponensial. \istilah{Variable elimination} memanfaatkan
faktorisasi: penjumlahan didorong sedalam mungkin ke dalam hasil kali, sehingga setiap variabel
dijumlahkan hanya dari faktor-faktor yang menyebutnya. Hasil antara disimpan sebagai faktor baru.

Urutan eliminasi menentukan biaya. Urutan yang baik menjaga faktor antara tetap kecil,
sedangkan urutan yang buruk bisa membuat faktor sebesar tabel gabungan. Untuk graf berbentuk
pohon, biayanya linear. Untuk graf umum, inferensi eksak termasuk masalah NP-hard, dan biayanya
bergantung secara eksponensial pada ukuran faktor terbesar yang tidak bisa dihindari.

Contoh rantai tiga variabel biner membuat penghematannya terlihat. Misalkan $A\to B\to C$ dengan $p(A{=}1)=0{,}3$,
$p(B{=}1\mid A)$ bernilai $0{,}2$ untuk $A=0$ dan $0{,}9$ untuk $A=1$, serta $p(C{=}1\mid B)$ bernilai $0{,}1$ untuk $B=0$ dan $0{,}8$ untuk
$B=1$. Untuk menghitung $p(C{=}1)=\sum_a\sum_bp(a)\,p(b\mid a)\,p(C{=}1\mid b)$, penjumlahan atas $a$ didorong ke dalam karena hanya dua
faktor pertama yang menyebut $A$. Hasilnya faktor baru $\tau(b)=\sum_ap(a)p(b\mid a)$, dengan
$\tau(1)=0{,}7\cdot0{,}2+0{,}3\cdot0{,}9=0{,}41$ dan $\tau(0)=0{,}59$. Lalu $p(C{=}1)=0{,}59\cdot0{,}1+0{,}41\cdot0{,}8=0{,}387$. Untuk rantai $n$
variabel biner, menjumlahkan tabel gabungan membutuhkan $2^{n-1}$ suku, sedangkan eliminasi hanya beberapa operasi per variabel. Dengan
tiga puluh variabel, selisihnya antara setengah miliar suku dan seratusan operasi.

Kalau banyak kueri harus dijawab pada jaringan yang sama, eliminasi yang diulang untuk setiap kueri membuang kerja. Kuliah membahas
\istilah{clique tree}: variabel-variabel dikelompokkan ke dalam klik yang disusun sebagai pohon, lalu pesan dikirim sekali ke arah akar dan
sekali kembali ke daun. Setelah dua sapuan itu, setiap klik menyimpan distribusi marginal variabel-variabelnya, sehingga semua kueri
tunggal bisa dijawab sekaligus. Inilah message passing dari \cref{ch:gdl} dalam bentuk probabilistik. Pada graf yang punya siklus,
pesan yang sama kadang tetap dikirim berulang-ulang sampai stabil, versi \emph{loopy} yang tidak dijamin benar tetapi sering cukup baik
dalam praktik.

\section{Inferensi dengan sampel}

Kalau inferensi eksak terlalu mahal, kita bisa memperkirakan peluang dengan menarik sampel. Cara
paling sederhana adalah \istilah{forward sampling}: tarik setiap variabel mengikuti urutan graf, dari
induk ke anak. Untuk menghitung peluang bersyarat, \istilah{rejection sampling} membuang sampel yang
tidak cocok dengan pengamatan. Cara ini boros sekali kalau pengamatannya jarang terjadi, karena
hampir semua sampel dibuang.

\istilah{Likelihood weighting} tidak membuang apa pun. Variabel yang diamati dipaksa bernilai sesuai
pengamatan, dan setiap sampel diberi bobot sebesar peluang pengamatan itu terjadi. Itulah
contoh \istilah{importance sampling}: sampel ditarik dari distribusi yang mudah, lalu dikoreksi dengan
bobot. Kelemahannya muncul ketika pengamatan berada di bagian bawah graf. Variabel di atasnya
ditarik tanpa memperhatikan pengamatan, sehingga banyak sampel mendapat bobot yang hampir nol.

Markov chain Monte Carlo mengambil jalan yang berbeda. Kita membangun sebuah rantai Markov yang
distribusi stasionernya adalah distribusi posterior yang dicari, lalu menjalankan rantai itu cukup
lama. \istilah{Gibbs sampling} memperbarui satu variabel pada satu waktu, menariknya dari distribusi
bersyarat terhadap semua variabel lain. Berkat Markov blanket, distribusi bersyarat itu hanya
melibatkan beberapa faktor lokal. Metropolis--Hastings \citep{metropolis1953,hastings1970} lebih umum:
usulkan langkah acak, lalu terima dengan peluang yang dipilih tepat agar rantai menyeimbangkan
distribusi target. Kedua metode ini lahir dari fisika statistik, untuk menghitung sifat gas dan
cairan, dan kembali ke fisika di \cref{ch:fiskom}.

Bayangkan menaksir bentuk sebuah ruangan gelap dengan berjalan acak di dalamnya, sambil lebih sering
melangkah ke tempat yang lebih ``terang''. Setelah cukup lama, jejak langkah Anda menggambarkan
seberapa terang setiap sudut ruangan. Langkah-langkah awal sebelum rantai mencapai keseimbangan
biasanya dibuang, dan sampel yang berurutan saling berkorelasi, sehingga jumlah sampel efektifnya
lebih kecil dari jumlah sampel mentah.

\section{Belajar parameter dan struktur}

Kalau struktur graf diketahui dan datanya lengkap, belajar parameter ternyata hanya menghitung.
Maximum likelihood untuk $p(x_i\mid\mathrm{pa}(x_i))$ adalah frekuensi relatif di data. Masalahnya
muncul untuk kombinasi yang jarang atau tidak pernah terlihat, yang langsung mendapat peluang nol.
Pendekatan Bayesian memberi prior Dirichlet pada setiap tabel peluang, dan efeknya setara dengan
menambahkan hitungan semu sebelum melihat data. Kalau sebuah kejadian muncul 3 kali dari 10, maximum
likelihood memberi $0{,}3$. Dengan satu hitungan semu untuk setiap dari dua kemungkinan, yaitu
Laplace smoothing, perkiraannya menjadi $4/12\approx0{,}33$. Pengaruh hitungan semu ini memudar ketika
data bertambah.

Kalau sebagian data hilang, atau ada variabel tersembunyi yang tidak pernah diamati, hitungan
sederhana tidak bisa dilakukan. Algoritme EM \citep{dempster1977} bergantian antara dua langkah. Langkah
E menghitung distribusi variabel yang hilang dengan parameter saat ini, sehingga menghasilkan hitungan
yang diharapkan. Langkah M memperbarui parameter seolah-olah hitungan yang diharapkan itu adalah
data lengkap. Setiap putaran tidak menurunkan likelihood. EM akan muncul lagi sebagai kerangka di
balik VAE dan model campuran Gauss di \cref{ch:mlat}.

Belajar struktur graf lebih sulit, karena jumlah graf yang mungkin tumbuh super-eksponensial
terhadap jumlah variabel. Kuliah membahas dua keluarga pendekatan. Pendekatan berbasis kendala
menjalankan uji independensi bersyarat pada data, lalu membangun graf yang sesuai dengan hasil uji
itu. Pendekatan berbasis skor mencari graf yang memaksimalkan sebuah skor, misalnya BIC,
\begin{equation}
\mathrm{BIC}=\log p(\mathcal D\mid\hat\theta,\mathcal G)-\frac{\log N}{2}\,\dim(\mathcal G),
\end{equation}
yang menyeimbangkan kecocokan dengan data dan jumlah parameter. Tanpa suku penalti, graf yang
terhubung penuh selalu menang, dan itulah overfitting dalam bentuk struktur. Pencariannya biasanya
heuristik, misalnya hill climbing dengan menambah, menghapus, atau membalik satu sisi pada satu waktu,
atau tabu search yang melarang kembali ke langkah yang baru saja diambil. Pendekatan ketiga, Bayesian
model averaging, tidak memilih satu graf, tetapi merata-ratakan prediksi atas banyak graf yang
diberi bobot posterior.

\section{Model yang berjalan dalam waktu}

Banyak sistem berubah dari waktu ke waktu. \istilah{Dynamic Bayesian network} memodelkannya dengan
mengulang struktur yang sama di setiap langkah waktu, dengan dua asumsi: keadaan sekarang hanya
bergantung pada keadaan sebelumnya (asumsi Markov), dan aturan transisinya tidak berubah dari waktu
ke waktu. Seluruh model cukup ditentukan oleh satu irisan waktu beserta sambungannya ke irisan
berikutnya, lalu bisa ``dibuka'' sepanjang yang dibutuhkan.

Ada empat tugas inferensi yang berulang di semua model temporal. \istilah{Filtering} menghitung
keadaan sekarang dari semua pengamatan sampai sekarang. \istilah{Prediction} menghitung keadaan masa
depan. \istilah{Smoothing} memperbaiki perkiraan keadaan masa lalu dengan pengamatan yang datang
belakangan. Dan \istilah{most likely explanation} mencari urutan keadaan yang paling mungkin
menjelaskan semua pengamatan.

\subsection{Hidden Markov model}

Hidden Markov model \citep{rabiner1989} adalah DBN dengan satu variabel keadaan diskret yang tersembunyi
dan satu pengamatan di setiap langkah. Ia ditentukan oleh matriks transisi $a_{ij}=p(s_t=j\mid s_{t-1}=i)$,
peluang emisi $b_j(o)=p(o_t=o\mid s_t=j)$, dan distribusi awal. Filtering dihitung dengan
\istilah{forward algorithm},
\begin{equation}
\alpha_t(j)=b_j(o_t)\sum_i\alpha_{t-1}(i)\,a_{ij},
\end{equation}
yang bergantian antara memprediksi (penjumlahan atas keadaan sebelumnya) dan memperbarui dengan
pengamatan baru (perkalian dengan peluang emisi). Algoritme backward menghitung pesan yang sama dari
arah sebaliknya, dan gabungan keduanya memberi smoothing. Algoritme Viterbi mengganti penjumlahan
dengan maksimum, $\delta_t(j)=b_j(o_t)\max_i\delta_{t-1}(i)\,a_{ij}$, sambil mencatat $i$ mana yang memberi maksimum,
lalu menelusuri catatan itu mundur dari akhir untuk mendapatkan urutan keadaan terbaik. Biayanya sebanding dengan
panjang deret kali kuadrat jumlah keadaan, bukan eksponensial terhadap panjang deret seperti kalau semua urutan dicoba, dan Baum--Welch adalah EM untuk mempelajari
parameter HMM dari pengamatan saja.

Contoh klasik dari \citet{russell2021} adalah penjaga yang tidak pernah keluar gedung dan hanya bisa
menebak cuaca dari apakah direkturnya datang membawa payung. Misalkan peluang hujan hari ini sama
dengan kemarin adalah $0{,}7$, peluang membawa payung saat hujan $0{,}9$, dan saat tidak hujan
$0{,}2$, dengan keyakinan awal $0{,}5$. Pada hari pertama direktur membawa payung. Prediksinya tetap
$0{,}5$, dan pembaruan memberi hujan sebanding dengan $0{,}9\times0{,}5=0{,}45$ dan tidak hujan
sebanding dengan $0{,}2\times0{,}5=0{,}10$. Setelah dinormalisasi, peluang hujan $0{,}818$. Pada hari
kedua prediksinya $0{,}818\times0{,}7+0{,}182\times0{,}3\approx0{,}627$, dan pengamatan payung kedua
menaikkannya menjadi sekitar $0{,}883$. Setiap hari, keyakinan mengalir lewat dua langkah yang sama:
ramalkan, lalu koreksi dengan bukti.

\subsection{Kalman filter}

Kalau keadaannya kontinu dan semua hubungan linear dengan noise Gauss, filtering bisa dikerjakan
secara eksak dengan Kalman filter \citep{kalman1960}. Modelnya
$\vx_t=\mA\vx_{t-1}+\bm w_t$ dan $\bm z_t=\bm H\vx_t+\bm v_t$, dengan noise proses
$\bm w_t\sim\N(\bm0,\bm Q)$ dan noise pengukuran $\bm v_t\sim\N(\bm0,\bm R)$. Karena Gauss tetap Gauss di
bawah transformasi linear, cukup melacak rata-rata dan kovarians:
\begin{equation}
\begin{aligned}
\hat\vx^-_t&=\mA\hat\vx_{t-1}, & \bm P^-_t&=\mA\bm P_{t-1}\mA^\top+\bm Q,\\
\bm K_t&=\bm P^-_t\bm H^\top\big(\bm H\bm P^-_t\bm H^\top+\bm R\big)^{-1}, &&\\
\hat\vx_t&=\hat\vx^-_t+\bm K_t\big(\bm z_t-\bm H\hat\vx^-_t\big), & \bm P_t&=(\bm I-\bm K_t\bm H)\bm P^-_t.
\end{aligned}
\end{equation}
Baris pertama adalah prediksi, dan baris terakhir adalah koreksi. Matriks $\bm K_t$, \istilah{Kalman gain},
menentukan seberapa besar kita mempercayai pengukuran dibanding prediksi.

Kalman filter seperti mendengarkan dua saksi. Saksi pertama adalah model, yang meramalkan posisi
dari posisi sebelumnya. Saksi kedua adalah sensor. Keduanya bisa keliru, dan kita mempercayai
masing-masing sesuai ketelitiannya. Dalam satu dimensi, kalau ketidakpastian prediksi punya varians
4 dan pengukuran punya varians 1, gain-nya $4/(4+1)=0{,}8$: hasil akhir bergerak 80 persen ke arah
pengukuran, dan variansnya turun menjadi $0{,}8$, lebih kecil dari keduanya. Menggabungkan dua sumber
yang tidak sempurna menghasilkan perkiraan yang lebih baik dari masing-masing.

Rumus satu dimensi ini bisa diturunkan langsung dengan teorema Bayes. Prediksi memberi prior
$\N(\hat x^-,P^-)$, dan pengukuran $z$ memberi likelihood $\N(z;x,R)$. Posterior sebanding dengan hasil kali
keduanya, dan eksponennya adalah
$-\tfrac12\big[(x-\hat x^-)^2/P^-+(z-x)^2/R\big]$. Jumlah dua kuadrat dalam $x$ tetap kuadrat dalam $x$, jadi
posteriornya Gauss. Mengumpulkan koefisien $x^2$ memberi presisi posterior $1/P=1/P^-+1/R$: presisi,
yaitu kebalikan varians, langsung dijumlahkan. Mengumpulkan koefisien $x$ memberi rata-rata posterior
$\hat x=P\big(\hat x^-/P^-+z/R\big)$, rata-rata berbobot dengan bobot sebanding presisi masing-masing.
Dengan sedikit aljabar, keduanya bisa ditulis ulang sebagai $\hat x=\hat x^-+K(z-\hat x^-)$ dan
$P=(1-K)P^-$ dengan $K=P^-/(P^-+R)$, persis bentuk Kalman. Untuk angka di atas, $1/P=1/4+1/1=1{,}25$,
sehingga $P=0{,}8$.

Untuk sistem yang tidak linear, Kalman filter diperluas dengan melinearkan model di sekitar perkiraan
saat ini (\istilah{extended Kalman filter}), atau diganti dengan \istilah{particle filter} yang mewakili
distribusi dengan sekumpulan sampel berbobot. Keduanya dipakai di kendaraan otonom untuk melacak
posisi dari GPS, odometri, dan lidar.

Particle filter adalah filter Bayes yang sama dengan Kalman filter, hanya distribusinya diwakili sampel, bukan Gauss. Setiap langkah
punya tiga bagian. Setiap partikel digerakkan menurut model dinamika, termasuk noise-nya. Setiap partikel diberi bobot sebanding dengan
likelihood pengukuran di posisinya. Lalu partikel ditarik ulang sesuai bobotnya, sehingga partikel di tempat yang tidak mungkin mati dan
partikel di tempat yang mungkin berlipat. Contoh kecil: tiga partikel di posisi 1, 2, dan 3 bergerak satu satuan ke 2, 3, dan 4. Sensor
mengukur $3{,}2$ dengan noise Gauss bersimpangan baku satu. Likelihood-nya sebanding dengan $e^{-(x-3{,}2)^2/2}$, yaitu $0{,}49$, $0{,}98$, dan
$0{,}73$, dan setelah dinormalisasi menjadi bobot $0{,}22$, $0{,}45$, dan $0{,}33$. Taksiran posisinya rata-rata berbobot, sekitar $3{,}11$.
Dalam praktik dipakai ratusan sampai ribuan partikel, dan karena tidak ada asumsi Gauss, distribusinya boleh bercabang, misalnya ketika
robot belum tahu di lorong mana ia berada, persis lokalisasi Markov di \cref{ch:robot}.

\section{Di luar kelas}

Model-model di bab ini bekerja di belakang banyak sistem yang kita temui di bagian-bagian berikutnya.
Kendaraan otonom melacak posisinya dengan filter Bayes dan melacak mobil di sekitarnya dengan Kalman
filter (\cref{ch:robot}). Pelacak gerak inersia menggabungkan giroskop dan akselerometer dengan versi
sederhana Kalman filter (\cref{ch:pervasif}). Filter Wiener dan estimator MMSE dalam pemrosesan sinyal
adalah saudara dekatnya (\cref{ch:sinyal}). Dalam bioinformatika, hidden Markov model dipakai untuk
mencari gen dan menjajarkan keluarga protein \citep{durbin1998}, dan algoritme EM menjadi dasar banyak
metode analisis ekspresi gen (\cref{ch:genomik}).

Di dunia simulasi, ide Kalman filter dikenal sebagai \istilah{data assimilation}. Prakiraan cuaca
modern bekerja dengan cara ini: model atmosfer meramalkan keadaan beberapa jam ke depan, lalu
ramalan itu dikoreksi dengan pengamatan dari satelit, balon, dan stasiun cuaca, dengan bobot sesuai
ketidakpastian masing-masing. Model atmosfer punya jutaan variabel, sehingga matriks kovarians
$\bm P$ tidak mungkin disimpan. Ensemble Kalman filter \citep{evensen1994} memperkirakannya dari
sekumpulan simulasi yang berjalan paralel dengan kondisi awal sedikit berbeda. Sebaran ensemble
menjadi ukuran ketidakpastian ramalan.

Gagasan yang sama berguna di skala laboratorium. Simulasi CFD sebuah tangki dan pengukuran
temperatur dari sensor di dalam tangki yang sama jarang cocok sempurna. Model turbulensi,
sifat fluida, dan syarat batas semuanya menyimpan galat. Daripada memilih salah satu, kita bisa
menarik simulasi secara perlahan ke arah pengukuran di titik-titik sensor, dengan kekuatan yang
mencerminkan seberapa dipercaya masing-masing. Teknik paling sederhana seperti ini disebut
\istilah{nudging}, dan ia adalah kerabat dekat Kalman filter dengan gain yang dipilih tangan.

Cara berpikir probabilistik juga membuka pintu ke masalah invers. Kalau kita ingin menyimpulkan
parameter fisika, misalnya konduktivitas atau syarat batas, dari pengukuran yang berisik, jawaban
yang jujur bukan satu angka, melainkan sebuah distribusi posterior. MCMC, inferensi variasional, dan
model generatif dari \cref{ch:mlat} adalah alat untuk menghitungnya.

\section{Perkembangan riset}

Banyak ilmu punya simulator yang teliti tetapi tidak punya likelihood yang bisa dituliskan: simulator tabrakan
partikel, model epidemi berbasis agen, atau model iklim. Bayes membutuhkan $p(\text{data}\mid\theta)$, sedangkan
simulator hanya bisa menghasilkan data untuk $\theta$ tertentu. \istilah{Simulation-based inference} \citep{cranmer2020sbi}
menjembatani keduanya. Jalankan simulator untuk banyak $\theta$ yang ditarik dari prior, lalu latih jaringan saraf untuk
mendekati posterior $p(\theta\mid\text{data})$ secara langsung, misalnya dengan normalizing flow dari \cref{ch:arsitektur}.
Setelah dilatih, posterior untuk pengamatan baru diperoleh hampir seketika, tanpa menjalankan MCMC lagi. Inferensi
semacam ini disebut \emph{amortized}: biaya besar dibayar sekali di awal, lalu dipakai berulang kali. Gagasan yang sama
dipakai di asimilasi data dan dalam mengkalibrasi simulasi fisika terhadap percobaan.

\sumberbab{Bab ini mengikuti kuliah Probabilistic Models (Bayesian network, inferensi eksak dan
aproksimasi, pembelajaran parameter dan struktur, DBN, HMM, Kalman filter) beserta latihannya. Slide
kuliah bersandar pada \citet{koller2009} dan \citet{russell2021}, yang juga menjadi bacaan pendamping
terbaik. Bagian Kalman filter dan particle filter juga muncul di kuliah kendaraan otonom.}
